% GENERATED FILE. Edit canonical lessons, then run npm run generate:books.
% canonical-source-hash: fnv1a64:a941941450276552
% canonical-lessons: GU-C439-dungar, GU-C439-sikhar, GU-C439-khadak, GU-C439-kakro, GU-C439-padghavu, GU-R439-review-433-439

\chapter{Hills and Rocks}
\label{ch:gu-hills-and-rocks}
\hlchaptermodality{439}

\begin{chapteropening}
\textbf{By the end of this chapter:} \emph{I can describe hills, peaks and rocks.}

Complete the last lesson of chapter 439: I can describe hills, peaks and rocks.
\end{chapteropening}

\section[\texorpdfstring{\gu{ડુંગર} (ḍungar)}{ડુંગર (ḍungar)}]{\gu{ડુંગર} (ḍungar) — a rocky hill; a hill}

\label{lesson:GU-C439-dungar}



\begin{quote}
Before the new one: say the Gujarati for hazy, then the Gujarati for to twinkle.
\end{quote}

\subsection*{You'll want to know: \gu{ડુંગર}}
\textbf{\gu{ડુંગર}} — \emph{ḍungar} — \textquotedblleft{}a rocky hill; a hill\textquotedblright{}.

\begin{quote}
\gu{ગામની} \gu{પાછળ} \gu{મોટો} \gu{ડુંગર} \gu{છે}.
\end{quote}

\begin{quote}
\emph{gāmnī pāchaḷ moṭo ḍungar che.}
\end{quote}

\begin{quote}
There is a big hill behind the village.
\end{quote}

A masculine noun.

\begin{cousinweb}
\textbf{\gu{ડુંગર}} (\emph{ḍungar}) goes back to a group of forms CDIAL reconstructs for \textquotedblleft{}a hill\textquotedblright{}, \emph{*ṭakka} and \emph{*ṭuṅga}; CDIAL lists the Gujarati word, \textquotedblleft{}a hillock\textquotedblright{} (CDIAL 5423).
\end{cousinweb}

\begin{grammarlens}[title={What you've built}]
One more word for getting around.
\end{grammarlens}

\subsection*{Your turn}
\begin{itemize}
\raggedright
  \item \emph{Say it:} \emph{ḍungar}
  \item \emph{Say it:} \emph{ḍungar}, once more
  \item \emph{Say it:} \emph{gāmnī pāchaḷ moṭo ḍungar che.} — There is a big hill behind the village.
  \item \emph{Recall:} say the Gujarati for hazy, then the Gujarati for to twinkle, then say \emph{ḍungar} again
\end{itemize}

\begin{tcolorbox}[breakable,colback=teal!4,colframe=teal!35!black,title={Before you move on}]
What does \gu{ડુંગર} mean? (\textquotedblleft{}A rocky hill; a hill\textquotedblright{}.) Say it once more.
\end{tcolorbox}
\section[\texorpdfstring{\gu{શિખર} (śikhar)}{શિખર (śikhar)}]{\gu{શિખર} (śikhar) — a summit; a peak}

\label{lesson:GU-C439-sikhar}



\begin{quote}
Before the new one: say the Gujarati for to twinkle, then the Gujarati for a rocky hill.
\end{quote}

\subsection*{You'll want to know: \gu{શિખર}}
\textbf{\gu{શિખર}} — \emph{śikhar} — \textquotedblleft{}a summit; a peak\textquotedblright{}.

\begin{quote}
\gu{પર્વતના} \gu{શિખર} \gu{પર} \gu{બરફ} \gu{છે}.
\end{quote}

\begin{quote}
\emph{parvatnā śikhar par baraph che.}
\end{quote}

\begin{quote}
There is snow on the summit of the mountain.
\end{quote}

A neuter noun.

\begin{cousinweb}
\textbf{\gu{શિખર}} (\emph{śikhar}) is a tatsama, Sanskrit \emph{śikhara}, \textquotedblleft{}a point, a peak\textquotedblright{} (Monier-Williams); it is also the spire of a temple.
\end{cousinweb}

\begin{grammarlens}[title={What you've built}]
One more word for getting around.
\end{grammarlens}

\subsection*{Your turn}
\begin{itemize}
\raggedright
  \item \emph{Say it:} \emph{śikhar}
  \item \emph{Say it:} \emph{śikhar}, once more
  \item \emph{Say it:} \emph{parvatnā śikhar par baraph che.} — There is snow on the summit of the mountain.
  \item \emph{Recall:} say the Gujarati for to twinkle, then the Gujarati for a rocky hill, then say \emph{śikhar} again
\end{itemize}

\begin{tcolorbox}[breakable,colback=teal!4,colframe=teal!35!black,title={Before you move on}]
What does \gu{શિખર} mean? (\textquotedblleft{}A summit; a peak\textquotedblright{}.) Say it once more.
\end{tcolorbox}
\section[\texorpdfstring{\gu{ખડક} (khaḍak)}{ખડક (khaḍak)}]{\gu{ખડક} (khaḍak) — a crag; a big rock}

\label{lesson:GU-C439-khadak}



\begin{quote}
Before the new one: say the Gujarati for a rocky hill, then the Gujarati for a summit.
\end{quote}

\subsection*{You'll want to know: \gu{ખડક}}
\textbf{\gu{ખડક}} — \emph{khaḍak} — \textquotedblleft{}a crag; a big rock\textquotedblright{}.

\begin{quote}
\gu{દરિયાના} \gu{કિનારે} \gu{મોટા} \gu{ખડક} \gu{છે}.
\end{quote}

\begin{quote}
\emph{dariyānā kināre moṭā khaḍak che.}
\end{quote}

\begin{quote}
There are big rocks on the sea shore.
\end{quote}

A masculine noun.

\begin{cousinweb}
\textbf{\gu{ખડક}} (\emph{khaḍak}) No dictionary we use records an older form, so its history is not settled.
\end{cousinweb}

\begin{grammarlens}[title={What you've built}]
One more word for getting around.
\end{grammarlens}

\subsection*{Your turn}
\begin{itemize}
\raggedright
  \item \emph{Say it:} \emph{khaḍak}
  \item \emph{Say it:} \emph{khaḍak}, once more
  \item \emph{Say it:} \emph{dariyānā kināre moṭā khaḍak che.} — There are big rocks on the sea shore.
  \item \emph{Recall:} say the Gujarati for a rocky hill, then the Gujarati for a summit, then say \emph{khaḍak} again
\end{itemize}

\begin{tcolorbox}[breakable,colback=teal!4,colframe=teal!35!black,title={Before you move on}]
What does \gu{ખડક} mean? (\textquotedblleft{}A crag; a big rock\textquotedblright{}.) Say it once more.
\end{tcolorbox}
\section[\texorpdfstring{\gu{કાંકરો} (kā̃kro)}{કાંકરો (kā̃kro)}]{\gu{કાંકરો} (kā̃kro) — a pebble}

\label{lesson:GU-C439-kakro}



\begin{quote}
Before the new one: say the Gujarati for a summit, then the Gujarati for a crag.
\end{quote}

\subsection*{You'll want to know: \gu{કાંકરો}}
\textbf{\gu{કાંકરો}} — \emph{kā̃kro} — \textquotedblleft{}a pebble\textquotedblright{}.

\begin{quote}
\gu{બાળકે} \gu{તળાવમાં} \gu{કાંકરો} \gu{ફેંક્યો}.
\end{quote}

\begin{quote}
\emph{bāḷke taḷāvmā̃ kā̃kro phẽkyo.}
\end{quote}

\begin{quote}
The child threw a pebble into the pond.
\end{quote}

A masculine noun.

\begin{cousinweb}
\textbf{\gu{કાંકરો}} (\emph{kā̃kro}) is a tadbhava of Sanskrit \emph{karkara}, \textquotedblleft{}a stone\textquotedblright{}; CDIAL lists the Gujarati word, \textquotedblleft{}a pebble\textquotedblright{} (CDIAL 2820).
\end{cousinweb}

\begin{grammarlens}[title={What you've built}]
One more word for getting around.
\end{grammarlens}

\subsection*{Your turn}
\begin{itemize}
\raggedright
  \item \emph{Say it:} \emph{kā̃kro}
  \item \emph{Say it:} \emph{kā̃kro}, once more
  \item \emph{Say it:} \emph{bāḷke taḷāvmā̃ kā̃kro phẽkyo.} — The child threw a pebble into the pond.
  \item \emph{Recall:} say the Gujarati for a summit, then the Gujarati for a crag, then say \emph{kā̃kro} again
\end{itemize}

\begin{tcolorbox}[breakable,colback=teal!4,colframe=teal!35!black,title={Before you move on}]
What does \gu{કાંકરો} mean? (\textquotedblleft{}A pebble\textquotedblright{}.) Say it once more.
\end{tcolorbox}
\section[\texorpdfstring{\gu{પડઘાવું} (paḍghāvũ)}{પડઘાવું (paḍghāvũ)}]{\gu{પડઘાવું} (paḍghāvũ) — to echo}

\label{lesson:GU-C439-padghavu}



\begin{quote}
Before the new one: say the Gujarati for a crag, then the Gujarati for a pebble.
\end{quote}

\subsection*{You'll want to know: \gu{પડઘાવું}}
\textbf{\gu{પડઘાવું}} — \emph{paḍghāvũ} — \textquotedblleft{}to echo\textquotedblright{}.

\begin{quote}
\gu{ખીણમાં} \gu{અમારો} \gu{અવાજ} \gu{પડઘાયો}.
\end{quote}

\begin{quote}
\emph{khīṇmā̃ amāro avāj paḍghāyo.}
\end{quote}

\begin{quote}
Our voices echoed in the valley.
\end{quote}

Past (it …): \gu{તે} \gu{અવાજ} \gu{પડઘાયો}, \emph{te avāj paḍghāyo}.

\begin{cousinweb}
\textbf{\gu{પડઘાવું}} (\emph{paḍghāvũ}) is built on \gu{પડઘો}, an echo. CDIAL lists the Kachchhi word \emph{paṛgho}, \textquotedblleft{}an echo\textquotedblright{}, under \emph{*pratighōṣa}, \textquotedblleft{}echo\textquotedblright{}, a form it reconstructs (CDIAL 8555a); it gives no Gujarati form, so the link is hedged.
\end{cousinweb}

\begin{grammarlens}[title={What you've built}]
One more word for getting around.
\end{grammarlens}

\subsection*{Your turn}
\begin{itemize}
\raggedright
  \item \emph{Say it:} \emph{paḍghāvũ}
  \item \emph{Say it:} \emph{paḍghāvũ}, once more
  \item \emph{Say it:} \emph{khīṇmā̃ amāro avāj paḍghāyo.} — Our voices echoed in the valley.
  \item \emph{Recall:} say the Gujarati for a crag, then the Gujarati for a pebble, then say \emph{paḍghāvũ} again
\end{itemize}

\begin{tcolorbox}[breakable,colback=teal!4,colframe=teal!35!black,title={Before you move on}]
What does \gu{પડઘાવું} mean? (\textquotedblleft{}To echo\textquotedblright{}.) Say it once more.
\end{tcolorbox}
\section[(dialogue)]{Looking back: from Outfits and Bags to Hills and Rocks}

\label{lesson:GU-R439-review-433-439}



\begin{quote}
Say the Gujarati for an outfit and a shoulder bag.
\end{quote}

\subsection*{Your turn}
Say these aloud:

\begin{itemize}
\raggedright
  \item an outfit, a shoulder bag, a sack, to suit, to shrink
  \item a sharp shower, a drizzle, a raincoat, to gather (of clouds), to get wet through
  \item a crash of thunder, a whirlwind, a howl of wind, a forecast, to flash
  \item a hot summer wind, steamy heat, mercury, to scorch, to bake (in the heat)
  \item a biting chill, heavy dew, the cold, to congeal, to shiver all over
  \item the light of daybreak, dusk, dim, hazy, to twinkle
  \item a rocky hill, a summit, a crag, a pebble, to echo
\end{itemize}

\begin{tcolorbox}[breakable,colback=teal!4,colframe=teal!35!black,title={Before you move on}]
An outfit? (\textbf{\gu{પોશાક}}, \emph{pośāk}.) A shoulder bag? (\textbf{\gu{ઝોળી}}, \emph{jhoḷī}.) A sack? (\textbf{\gu{કોથળો}}, \emph{kothḷo}.) To suit? (\textbf{\gu{શોભવું}}, \emph{śobhvũ}.) To shrink? (\textbf{\gu{સંકોચાવું}}, \emph{sankocāvũ}.) A sharp shower? (\textbf{\gu{ઝાપટું}}, \emph{jhāpṭũ}.) A drizzle? (\textbf{\gu{ઝરમર}}, \emph{jharmar}.) A raincoat? (\textbf{\gu{રેઈનકોટ}}, \emph{reīnkoṭ}.) To gather (of clouds)? (\textbf{\gu{ઘેરાવું}}, \emph{gherāvũ}.) To get wet through? (\textbf{\gu{ભીંજાવું}}, \emph{bhī̃jāvũ}.) A crash of thunder? (\textbf{\gu{કડાકો}}, \emph{kaḍāko}.) A whirlwind? (\textbf{\gu{ચક્રવાત}}, \emph{cakravāt}.) A howl of wind? (\textbf{\gu{સુસવાટો}}, \emph{susvāṭo}.) A forecast? (\textbf{\gu{આગાહી}}, \emph{āgāhī}.) To flash? (\textbf{\gu{ઝબકવું}}, \emph{jhabakvũ}.) A hot summer wind? (\textbf{\gu{લૂ}}, \emph{lū}.) Steamy heat? (\textbf{\gu{બફારો}}, \emph{baphāro}.) Mercury? (\textbf{\gu{પારો}}, \emph{pāro}.) To scorch? (\textbf{\gu{દઝાડવું}}, \emph{dajhāḍvũ}.) To bake (in the heat)? (\textbf{\gu{શેકાવું}}, \emph{śekāvũ}.) A biting chill? (\textbf{\gu{ઠાર}}, \emph{ṭhār}.) Heavy dew? (\textbf{\gu{ઝાકળ}}, \emph{jhākaḷ}.) The cold? (\textbf{\gu{ટાઢ}}, \emph{ṭāḍh}.) To congeal? (\textbf{\gu{થીજવું}}, \emph{thījvũ}.) To shiver all over? (\textbf{\gu{થરથરવું}}, \emph{thartharvũ}.) The light of daybreak? (\textbf{\gu{ઉજાસ}}, \emph{ujās}.) Dusk? (\textbf{\gu{સંધ્યા}}, \emph{sandhyā}.) Dim? (\textbf{\gu{ઝાંખું}}, \emph{jhā̃khũ}.) Hazy? (\textbf{\gu{ધૂંધળું}}, \emph{dhū̃dhḷũ}.) To twinkle? (\textbf{\gu{ટમટમવું}}, \emph{ṭamṭamvũ}.) A rocky hill? (\textbf{\gu{ડુંગર}}, \emph{ḍungar}.) A summit? (\textbf{\gu{શિખર}}, \emph{śikhar}.) A crag? (\textbf{\gu{ખડક}}, \emph{khaḍak}.) A pebble? (\textbf{\gu{કાંકરો}}, \emph{kā̃kro}.) To echo? (\textbf{\gu{પડઘાવું}}, \emph{paḍghāvũ}.)
\end{tcolorbox}
